\chapter{实证结果与分析}
\noindent\textbf{本章状态：}本版尚未运行公司数据的正式市场实验。以下内容给出结果组织、表格位置和对应论证问题；破折号表示未生成结果，不表示零损失或无显著差异。正式运行后，应以结果文件替换占位内容，并据真实证据完成段落。

\section{样本与风险网络的描述性结果}
\subsection{数据覆盖及主要变量}
本节首先报告最终样本形成日、证券构成、有效区间和各处理阶段的样本数量，使读者能够看清从原始分钟记录到未来方差标签的变化。描述性统计包括RV、参考价收益、条件sigma、CoVaR强度和有效边比例，并将缺失、停牌、未成熟标签及模型失败分别列示。

\Pending{表5.1：逐阶段样本损失与变量统计；图5.1：逐股逐日覆盖；图5.2：训练、验证、测试及成熟边界。}

风险网络描述先比较原始和标准化强度，再讨论条件输入、输出差异。网络图只辅助展示若干预先确定状态，固定节点位置与边尺度；不能用布局变化制造结构变化的视觉印象。

\subsection{风险定义与代数检查}
本节列示上下尾镜像误差、同形边际下的标准化双向差、不同边际自由度的分解及图支持非对称率。这些诊断用于确认实现符合第三章的数学设定，不把标准t-Copula的结构对称性当作数据中“发现没有非对称风险”的结论。

\Pending{表5.2：数值与边语义诊断；图5.3：原始方向差与标准化方向差。}

\section{整体样本外预测表现}
表\ref{tab:main_results_pending}保留全部主要模型与期限。正文应首先报告相对于强基线的差值和区间，再说明最优均值排序是否稳定；不能只展示表现最好的种子或删去主模型没有获胜的期限。

\begin{table}[htbp]\centering\small
\caption{主预测结果表结构：QLIKE（尚未生成实证值）}\label{tab:main_results_pending}
\begin{tabularx}{\textwidth}{Yccc}
\toprule 模型 & 1日 & 5日 & 22日\\\midrule
传统HAR / HARQ & — & — & —\\
匹配自身MLP & — & — & —\\
公共信息增强模型 & — & — & —\\
匹配风险汇总MLP & — & — & —\\
正则化横截面HAR & — & — & —\\
GLASSO图模型 & — & — & —\\
DCC相关边模型 & — & — & —\\
CoVaR标量边模型 & — & — & —\\
底层参数边模型 & — & — & —\\
完整边特征模型 & — & — & —\\\bottomrule
\end{tabularx}
\end{table}

正式表格同时标注训练损失、节点信息组、有效股票日及日期数。MSE、RMSE、MAE另表给出，不把不同单位的改善画在同一数值轴上。可用相对损失便于阅读，但仍保留原始损失及相应分母。

\section{预测改善来自何处}
\subsection{训练损失与公共信息}
首先比较同一模型从MSE训练到QLIKE训练的变化，再比较同一QLIKE训练下无图、公共信息增强和风险汇总模型。这样可判断总体改善中有多少在使用具体图结构之前已经出现。

\Pending{表5.4：训练损失交叉比较；图5.4：损失准则与公共信息贡献拆分。}

若完整模型仅超过MSE训练的HAR，但未超过同损失强基线，结论应限于整体流程或训练方式的改善，不将全部差异记为逐边信息贡献。

\subsection{风险汇总与逐边条件化}
核心比较为\code{M01_FULL}与\code{B06_MATCHED}。两者共享节点信息、风险汇总、自身MLP和训练方案，差别在逐边分支。应报告平均损失差、相对降低、日期区块置信区间以及不同年份的稳定性，并检查预测失败或缺边是否改变有效域。

\Pending{表5.5：核心成对比较与多重校正；图5.5：相对匹配汇总基线的累计日期损失差。}

若逐边分支的改进得到支持，说明具体关系在研究样本中保留了汇总变量之外可用于预测的结构；若区间较宽，则只能说明目前证据尚不足以精确区分，而不将点估计排序当作稳定优越性。

\subsection{拓扑、标量与多维边}
在固定支持下比较二值边、标量边和向量边，进而删除标准化、反向和变化通道。对每个消融报告重训还是推理屏蔽、参数量与入度是否变化。关系重配与随机重连结果用于检查“边值属于哪个邻居”是否重要，不将这些计算实验称为真实市场因果干预。

\Pending{表5.6：边特征与关系破坏消融；图5.6：成对差值及区间。}

\subsection{边际尺度与风险表征}
比较标准化风险边、动态相关边以及包含当前和滞后底层参数的模型。若风险表征优于完整参数对照，讨论其有限样本表示效率；若收益被sigma或公共风险解释，则按相应层次限定贡献。上下尾或反向通道的代数重复，不应通过重复检验获得多项独立贡献的表象。

\section{期限、市场状态与稳健性}
本节按预设的1、5、22日目标和预测时可知状态报告结果。状态表需包含实际日期数量和成对差异的组间检验。年份分析与累计损失差共同用于检查结果是否由少数日期驱动，不只展示模型最有利的危机区间。

稳健性表覆盖5/10/15分钟采样、500/750/1000日风险窗口、不同入边数及点/事件条件定义。每项结果附实际有效样本；若改变测度或窗口导致评价域变化，应再提供共同域比较。计算成本单独报告，不将更多算力投入隐藏在方法名称之后。

\Pending{表5.7：期限与状态；表5.8：稳健性与成本。}

\section{本章小结的回填原则}
正式小结应依次回答：是否超越强自身基线、是否超越公共信息和风险汇总、哪些边内容有独立增量、效果在哪些期限与状态成立、复杂度代价是多少。每项结论对应明确结果表，不以方法名称或图示替代证据。本稿不预先填写这些市场结论。
